% (a) 资源价格（对偶）可行域 2u1+u2>=3, u1+2u2>=2, u>=0 与等值线 4u1+4u2=k（k=20/3,8,10），最优价格(4/3,1/3)。
% (b) 数轴：原始可行点给出收益 0,5,6（下界），对偶可行价格(1,1),(2,1)给出 8,10（上界），二者在20/3相遇。
\begin{tikzpicture}[note/.style={font=\FigureNoteFont\fontsize{7.5}{9}\selectfont,text=FigureMuted},
  ttl/.style={font=\fontsize{8}{10}\selectfont,anchor=north}]
  \begin{scope}[x=11mm,y=11mm]
    \fill[FigureGreenFill] (0,3)--(4/3,1/3)--(2,0)--(3.3,0)--(3.3,3.4)--(0,3.4)--cycle;
    \draw[FigureGreen,line width=.9pt] (0,3.4)--(0,3)--(4/3,1/3)--(2,0)--(3.3,0);
    \begin{scope}\clip (0,0) rectangle (3.3,3.4);
      \foreach \k in {2,2.5}{\draw[FigureBlue,line width=.6pt,dashed] (0,\k)--(\k,0);}
      \draw[FigureBlue,line width=1.1pt] (0,5/3)--(5/3,0);
    \end{scope}
    \draw[->,FigureStroke,line width=.6pt] (-0.15,0)--(3.45,0) node[right]{$u_1$};
    \draw[->,FigureStroke,line width=.6pt] (0,-0.15)--(0,3.55) node[above]{$u_2$};
    \fill[FigureRed] (4/3,1/3) circle (1.7pt);
    \fill[FigureInk] (1,1) circle (1.3pt);
    \fill[FigureInk] (2,1) circle (1.3pt);
    \node[note,anchor=south west] at (1.02,1.02) {$8$};
    \node[note,anchor=south west] at (2.02,1.02) {$10$};
    \draw[FigureMuted,line width=.4pt] (1.36,0.29)--(1.5,-0.08);
    \node[note,anchor=north west] at (1.4,-0.04) {$(4/3,1/3)$，$20/3$};
    \draw[->,FigureRed,line width=.9pt] (2.6,2.6)--++(-0.45,-0.45);
    \node[note,anchor=south west] at (2.55,2.6) {降低总价};
    \node[note] at (2.1,1.75) {价格可行域};
    \node[ttl] at (1.6,-0.42) {(a) 资源价格的可行域};
  \end{scope}
  \begin{scope}[xshift=48mm,x=5.6mm,y=10mm]
    \draw[->,FigureStroke,line width=.6pt] (-0.3,0)--(11,0);
    \foreach \t in {0,2,4,6,8,10}{\draw[FigureStroke,line width=.5pt] (\t,0)--(\t,-0.08) node[below,font=\fontsize{7.5}{9}\selectfont]{$\t$};}
    \foreach \v in {0,5,6}{\fill[FigureBlue] (\v,0.35) circle (1.6pt);}
    \foreach \v in {8,10}{\fill[FigureGreen] (\v,0.95) circle (1.6pt);}
    \draw[FigureRed,line width=1.1pt] (20/3,-0.05)--(20/3,1.35);
    \node[note,anchor=south,text=FigureRed!80!FigureInk] at (20/3,1.35) {$20/3$};
    \node[note,anchor=south west,text=FigureBlue!80!FigureInk] at (-0.3,0.45) {原始可行获利};
    \node[note,anchor=south east,text=FigureGreen!75!FigureInk] at (10.6,1.05) {对偶可行上界};
    \draw[FigureMuted,line width=.5pt,|-|] (6,2.0)--(8,2.0);
    \node[note,anchor=south] at (7,2.03) {间隙$\le2$};
    \node[ttl] at (5.3,-0.55) {(b) 两侧的界};
  \end{scope}
\end{tikzpicture}
